\textbf{(a) Modelo matemático:}

O crescimento exponencial de uma colônia bacteriana é descrito por:
\[
N(t)=N_0\,e^{kt},
\]
onde $N_0$ é a população inicial ($t=0$), $k>0$ é a taxa de crescimento e $t$ é o tempo.

\textbf{(b) Linearização — método dos mínimos quadrados em escala log:}

Tomando o logaritmo natural:
\[
\ln N(t) = \ln N_0 + k\,t.
\]
Definindo $y_i = \ln N_i$ e $x_i = t_i$, temos a regressão linear $y=a+bx$ com:
\[
b = k = \frac{n\sum x_i y_i - (\sum x_i)(\sum y_i)}{n\sum x_i^2 - (\sum x_i)^2},
\]
\[
a = \ln N_0 = \frac{\sum y_i - b\sum x_i}{n}, \quad N_0=e^a.
\]

\textbf{Exemplo com dados típicos de colônia bacteriana}
(com tempo de duplicação $\approx 20\,\mathrm{min}$):

\begin{center}
\begin{tabular}{ccc}
\hline
$t$ (min) & $N$ (UFC) & $\ln N$ \\
\hline
0 & 1000 & 6{,}908 \\
20 & 2000 & 7{,}601 \\
40 & 4000 & 8{,}294 \\
60 & 8000 & 8{,}987 \\
80 & 16000 & 9{,}680 \\
\hline
\end{tabular}
\end{center}

Calculando as somas ($n=5$):
\[
\sum t_i = 200,\quad \sum y_i = 41{,}470,\quad \sum t_i y_i = 1763{,}4,\quad \sum t_i^2=12000.
\]

\[
b = k = \frac{5\times1763{,}4 - 200\times41{,}470}{5\times12000 - 200^2}
= \frac{523}{20000}\approx0{,}03466\,\mathrm{min}^{-1}.
\]
\[
a=\frac{41{,}470 - 0{,}03466\times200}{5}=\frac{41{,}470-6{,}931}{5}=\frac{34{,}539}{5}\approx6{,}908.
\]
\[
N_0=e^{6{,}908}\approx1000\,\text{UFC}. \quad\checkmark
\]

O tempo de duplicação é:
\[
t_{\text{dup}}=\frac{\ln 2}{k}=\frac{0{,}6931}{0{,}03466}\approx20{,}0\,\mathrm{min}.\quad\checkmark
\]

\textbf{Equação final:}
\[
\boxed{N(t)=1000\,e^{0{,}0347\,t},}
\]
com $t$ em minutos e $N$ em unidades formadoras de colônia.

